\chapter{Nijenhuis operators, complex and Kähler manifolds}
\label{chap:complex}

Complex analytic manifolds are introduced in Chapter~\ref{chap:analytic}. This chapter studies
fields of endomorphisms of the tangent bundle, their Nijenhuis torsion, almost complex structures
and Kähler structures. As for metrics and forms, fields of endomorphisms are stored in the
coordinates of the preferred chart at each point, and their smoothness is that of their local
representatives. We follow the conventions of \cite{GLT}.

\begin{definition}[Fields of endomorphisms]
  \label{def:end-field}
  \lean{LocallyConvex.Manifolds.jLocalRep,
    LocallyConvex.Manifolds.EndomorphismField,
    LocallyConvex.Manifolds.EndomorphismField.loc,
    LocallyConvex.Manifolds.EndomorphismField.loc_eq,
    LocallyConvex.Manifolds.EndomorphismField.loc_add,
    LocallyConvex.Manifolds.EndomorphismField.loc_smul,
    LocallyConvex.Manifolds.EndomorphismField.loc_pref,
    LocallyConvex.Manifolds.EndomorphismField.smoothMBOn_loc,
    LocallyConvex.Manifolds.EndomorphismField.applyField,
    LocallyConvex.Manifolds.EndomorphismField.isVectorField_applyField,
    LocallyConvex.Manifolds.EndomorphismField.localRep_applyField}
  \leanok
  \uses{def:mb-manifold, def:tangent-bundle}
  A \emph{smooth field of endomorphisms} $N$ of $TM$ is a field of real linear maps $N_x$ of the
  tangent spaces whose local representatives $(y, v) \mapsto N_i(y) v$ are smooth. It acts on
  vector fields, $(NX)(x) = N_x X(x)$, and the local representative of $NX$ is $N_i X_i$.
\end{definition}

\begin{definition}[The Nijenhuis torsion]
  \label{def:nijenhuis}
  \lean{LocallyConvex.Manifolds.endD, LocallyConvex.Manifolds.nijenhuisLoc,
    LocallyConvex.Manifolds.endApply,
    LocallyConvex.Manifolds.EndomorphismField.nijenhuisAt}
  \leanok
  \uses{def:end-field}
  The \emph{Nijenhuis torsion} of $N$ is
  $\Omega_N(X, Y) = N[NX, Y] + N[X, NY] - [NX, NY] - N^2[X, Y]$ \cite[Definition 3.1]{GLT}. In
  a chart it is
  $\Omega_N(u, v) = N\big((\partial_u N) v - (\partial_v N) u\big) + (\partial_{Nv} N) u -
  (\partial_{Nu} N) v$.
\end{definition}

\begin{theorem}[The Nijenhuis torsion is a tensor]
  \label{thm:nijenhuis-bracket}
  \lean{LocallyConvex.Manifolds.localBracket_nijenhuis,
    LocallyConvex.Manifolds.localBracket_congr,
    LocallyConvex.Manifolds.EndomorphismField.lieBracket_nijenhuis}
  \leanok
  \uses{def:nijenhuis, lem:man-local-bracket}
  For $C^1$ local vector fields $X$, $Y$,
  $N[NX, Y] + N[X, NY] - [NX, NY] - N^2[X, Y] = \Omega_N(X, Y)$, with $\Omega_N$ given by the
  local formula of Definition~\ref{def:nijenhuis}. In particular, for $C^1$ vector fields on $M$,
  the torsion at $x$ only depends on $X(x)$ and $Y(x)$ \cite[Theorem 3.2]{GLT}.
\end{theorem}

\begin{proof}
  \leanok
  Expand the brackets with the product rule $\dd(NX)(w) = (\partial_w N) X + N\,\dd X(w)$: all
  the derivatives of $X$ and $Y$ cancel. For vector fields on $M$, compute in the preferred chart
  at $x$, where the local representative of $NX$ is $N_i X_i$.
\end{proof}

\begin{definition}[Almost complex structures]
  \label{def:almost-complex}
  \lean{LocallyConvex.Manifolds.AlmostComplexStructure,
    LocallyConvex.Manifolds.AlmostComplexStructure.loc_sq,
    LocallyConvex.Manifolds.AlmostComplexStructure.IsIntegrable,
    LocallyConvex.Manifolds.AlmostComplexStructure.IsIntegrable.lieBracket}
  \leanok
  \uses{def:end-field, thm:nijenhuis-bracket}
  An \emph{almost complex structure} is a smooth field of endomorphisms $J$ with $J^2 = -1$. It
  is \emph{integrable} if its Nijenhuis torsion vanishes in every chart; then
  $J[JX, Y] + J[X, JY] - [JX, JY] + [X, Y] = 0$ for $C^1$ vector fields.
\end{definition}

\begin{theorem}[The almost complex structure of a complex manifold]
  \label{thm:canonical-J}
  \lean{LocallyConvex.Manifolds.tangentCoordChange_complex_smul,
    LocallyConvex.Manifolds.mulI, LocallyConvex.Manifolds.mulI_apply,
    LocallyConvex.Manifolds.mulI_mulI, LocallyConvex.Manifolds.canonicalJ,
    LocallyConvex.Manifolds.canonicalJ_loc,
    LocallyConvex.Manifolds.canonicalJ_isIntegrable}
  \leanok
  \uses{def:almost-complex}
  On a complex analytic manifold, multiplication by $i$ in the charts is a well defined almost
  complex structure, and it is integrable.
\end{theorem}

\begin{proof}
  \leanok
  The transition maps are complex analytic, so their derivatives are complex linear
  (\cite{LCA}, \texttt{complexAnalyticOn\_iff\_holomorphicOn}); hence the local representatives
  of $J$ are multiplication by $i$ in every chart. Being constant, they have vanishing Nijenhuis
  torsion.
\end{proof}

The converse is the Newlander--Nirenberg theorem \cite{NN57}. In infinite dimensions it holds for
real analytic almost complex structures on real analytic Banach manifolds
\cite[Theorem 7, p.~64]{Bel05} (Theorem~\ref{thm:nn-manifold} below), but it fails for smooth structures on some Banach manifolds
\cite[Theorem 3.1, p.~401]{Pat99}; for Fréchet manifolds, Lempert points out loop spaces with integrable
almost complex structures that are not locally biholomorphic to locally convex spaces
\cite[\S2, pp.~489--490]{Lem98}.

\begin{definition}[Real analytic structures and holomorphic charts]
  \label{def:holomorphic-chart}
  \lean{LocallyConvex.Manifolds.AlmostComplexStructure.IsRealAnalytic,
    LocallyConvex.Manifolds.AlmostComplexStructure.IsHolomorphicChart}
  \leanok
  An almost complex structure $J$ is real analytic if its local representative in every chart is
  a real analytic map into $L(E)$. A $J$-holomorphic chart is a chart $\psi$ into a complex normed
  space, smooth with smooth inverse in the real sense, with $D\psi(J v) = i\, D\psi(v)$.
\end{definition}

\begin{theorem}[Newlander--Nirenberg, real analytic case]
  \label{thm:nn-manifold}
  \lean{LocallyConvex.Manifolds.AlmostComplexStructure.exists_holomorphic_chart,
    LocallyConvex.Manifolds.AlmostComplexStructure.IsHolomorphicChart.complexAnalyticOn_transition,
    LocallyConvex.Manifolds.nijenhuis_eq_neg_nijenhuisLoc}
  \leanok
  \uses{def:holomorphic-chart}
  Let $M$ be modeled on a real Banach space and $J$ an integrable real analytic almost complex
  structure on $M$. Then every point lies in the domain of a $J$-holomorphic chart into
  $\ker(J(x_0)_{\mathbb{C}} - i)$, and the transition maps between $J$-holomorphic charts are
  complex analytic \cite[Th\'eor\`eme 3.5]{Pen70}.
\end{theorem}

\begin{proof}
  \leanok
  In the preferred chart at $x_0$, $J$ is a real analytic $\hat J$ with $\hat J^2 = -1$ and
  vanishing Nijenhuis tensor; the local theorem of \cite{LCA} gives a chart $\varphi$ with
  $D\varphi \circ \hat J = i\,D\varphi$, and $\psi = \varphi \circ \mathrm{chart}$ is
  $J$-holomorphic. For two such charts, the real derivative of $\psi_2 \circ \psi_1^{-1}$ is
  $B \circ A$, where $A$ intertwines $i$ and $J$ and $B$ intertwines $J$ and $i$; so it is complex
  linear, and the transition map is complex differentiable, hence complex analytic.
\end{proof}

\begin{definition}[Hermitian structures and the fundamental form]
  \label{def:hermitian}
  \lean{LocallyConvex.Manifolds.altOfBilin, LocallyConvex.Manifolds.HermitianStructure,
    LocallyConvex.Manifolds.HermitianStructure.g_J_self,
    LocallyConvex.Manifolds.HermitianStructure.fundamentalForm,
    LocallyConvex.Manifolds.HermitianStructure.fundamentalForm_apply}
  \leanok
  \uses{def:almost-complex, def:diff-form}
  A \emph{Hermitian structure} is an almost complex structure $J$ together with a smooth field of
  symmetric bilinear forms $g$, not necessarily positive, with $g(Ju, Jv) = g(u, v)$. Its
  \emph{fundamental form} is the $2$-form $\omega(u, v) = g(Ju, v)$.
\end{definition}

\begin{theorem}[Kähler structures are symplectic]
  \label{thm:kaehler-symplectic}
  \lean{LocallyConvex.Manifolds.KaehlerStructure,
    LocallyConvex.Manifolds.KaehlerStructure.symplecticForm}
  \leanok
  \uses{def:hermitian, def:symplectic}
  A \emph{Kähler structure} is a Hermitian structure with $J$ integrable, $g$ weakly
  nondegenerate and $\omega$ closed. Its fundamental form is a weak symplectic form.
\end{theorem}

\begin{proof}
  \leanok
  If $\omega(v, \cdot) = g(Jv, \cdot) = 0$, then $Jv = 0$ by weak nondegeneracy of $g$, so
  $v = -J^2 v = 0$. Integrability is not used.
\end{proof}

\begin{example}[The flat Kähler structure]
  \label{ex:flat-kaehler}
  \lean{LocallyConvex.Manifolds.reInner, LocallyConvex.Manifolds.reInner_apply,
    LocallyConvex.Manifolds.continuous_reInner, LocallyConvex.Manifolds.flatMetric,
    LocallyConvex.Manifolds.flatHermitian, LocallyConvex.Manifolds.flatHermitian_isClosed,
    LocallyConvex.Manifolds.flatKaehler, LocallyConvex.Manifolds.flatKaehler_form_apply}
  \leanok
  \uses{thm:kaehler-symplectic, thm:canonical-J}
  A complex inner product space $E$, not necessarily complete, carries the Kähler structure
  $J = i$, $g = \operatorname{Re}\langle \cdot, \cdot \rangle$, with fundamental form
  $\omega(u, v) = \operatorname{Im}\langle u, v \rangle$.
\end{example}
